\section{Simulation Results and Empirical Validation}
\label{sec:experiments}

\begin{sourceblock}{Paragraph 1 - Section Overview}
In this section, we present the empirical results of our simulations to validate the theoretical findings established in Sections~\ref{sec:regression} and~\ref{sec:classification}. We begin by outlining the experimental setup, detailing the generative models for continuous and discrete prediction spaces, and then present the outcome of our parameter sweeps, offline learning rates, and online regret convergence.
\end{sourceblock}

\begin{faqblock}
\faqitem{Context}{What is the primary objective of these simulation experiments?}
\faqitem{Process}{What three aspects of the strategic model are evaluated?}
\faqanswers
\faqans{Context}{To empirically validate the theoretical results established for regression and classification.}
\faqans{Process}{(1) Parameter sweeps (decision sweeps); (2) Offline learning rates; (3) Online regret convergence.}
\end{faqblock}

\subsection{Experimental Setup}

\begin{sourceblock}{Paragraph 2 - Setup Overview}
To simulate realistic settings with asymmetric information and non-linear dependencies, we implement two distinct generative models:
\end{sourceblock}

\begin{faqblock}
\faqitem{Process}{Why are two distinct generative models implemented?}
\faqanswers
\faqans{Process}{To simulate settings with asymmetric information, non-linear dependencies, and both continuous/discrete prediction spaces.}
\end{faqblock}

\subsubsection{Continuous Space (Regression Model)}

\begin{sourceblock}{Paragraph 3 - Continuous Space Generative Model}
We generate the continuous variables $(X, Y, U) \in \mathbb{R}^3$ using a two-component Gaussian Mixture Model (GMM). This ensures that the data distribution is non-Gaussian, highlighting the generality of our framework:
\begin{enumerate}
    \item A latent component variable $z \in \{1, 2\}$ is drawn with equal probability.
    \item Conditioned on $z$, the variables are drawn from a multivariate Gaussian with mean vector $\mu_z$ and covariance matrix $\Sigma_z$. The means are shifted to create asymmetric regions: $\mu_1 = [1, 1, 1]^\top$ and $\mu_2 = [-1, -1, -1]^\top$.
    \item The covariance matrix is parameterized by the correlation coefficients $\rho_{XY}$ (correlation between features and labels) and $\rho_{XU}$ (correlation between features and sender preferences). The alignment between $Y$ and $U$ (beyond the features $X$) is controlled by a parameter $\gamma \in [0, 1]$, set to $\gamma = 0.8$ throughout. The standard Gaussians are combined as:
    \begin{align}
        X &= \mu_{z, X} + Z_X, \\
        Y &= \mu_{z, Y} + \rho_{XY} Z_X + \sqrt{1 - \rho_{XY}^2} Z_Y, \\
        U &= \mu_{z, U} + \rho_{XU} Z_X + \gamma \sqrt{1 - \rho_{XU}^2} Z_Y + \sqrt{1 - \gamma^2} \sqrt{1 - \rho_{XU}^2} Z_U,
    \end{align}
    where $Z_X, Z_Y, Z_U \sim \mathcal{N}(0, 1)$ are independent standard normal variables.
\end{enumerate}
\end{sourceblock}

\begin{faqblock}
\faqitem{Context}{Why is a Gaussian Mixture Model (GMM) used for regression?}
\faqitem{Process}{How are the correlation coefficients and alignment parameter configured?}
\faqanswers
\faqans{Context}{To generate non-Gaussian distributions, validating the generality of the continuous framework.}
\faqans{Process}{Covariance matrix uses coefficients $\rho_{XY}, \rho_{XU}$ and alignment $\gamma = 0.8$, combining independent standard normals $Z_X, Z_Y, Z_U$.}
\end{faqblock}

\subsubsection{Discrete Space (Classification Model)}

\begin{sourceblock}{Paragraph 4 - Discrete Space Generative Model}
For multiclass classification with $K$ classes, the preference $U \in \{1, \dots, K\}$ is drawn from a prior categorical distribution $\pi = [0.35, 0.4, 0.25]$ (for $K=3$). Conditioned on the preference $U = i$:
\begin{enumerate}
    \item The features $X \in \mathbb{R}^2$ are drawn from a class-conditional Gaussian:
    \begin{equation}
        X \mid U = i \sim \mathcal{N}(\mu_i, \sigma^2 I),
    \end{equation}
    where $\mu_1 = [-1, -1]^\top$, $\mu_2 = [0, 1.5]^\top$, and $\mu_3 = [1, -1]^\top$ represent centroids distributed on a circle, and the class variance is $\sigma^2 = 0.5$.
    \item The true label $Y \in \{1, \dots, K\}$ is drawn such that the sender's preference is aligned with $Y$ with probability $1 - \alpha$, and is uniformly distributed over the remaining $K-1$ classes with probability $\alpha$ (alignment noise):
    \begin{equation}
        \Pr(Y = j \mid U = i) = (1 - \alpha) \I{j = i} + \frac{\alpha}{K - 1} \I{j \ne i}.
    \end{equation}
\end{enumerate}
\end{sourceblock}

\begin{faqblock}
\faqitem{Context}{What centroids and variance parameters define the classification model?}
\faqitem{Process}{How is alignment noise $\alpha$ incorporated into the discrete label setup?}
\faqanswers
\faqans{Context}{Centroids $\mu_1 = [-1, -1]^\top$, $\mu_2 = [0, 1.5]^\top$, $\mu_3 = [1, -1]^\top$ and class variance $\sigma^2 = 0.5$.}
\faqans{Process}{Determines label alignment probability: $1-\alpha$ for matching sender preference, and $\alpha/(K-1)$ for other classes.}
\end{faqblock}

\subsection{Decision Sweeps}
\label{sec:exp_decision}

\begin{sourceblock}{Paragraph 5 - Decision Sweeps Overview and Fig 1}
We first validate our game-theoretic foundations by sweeping the environment parameters and plotting the resulting strategic advantage $\Delta^* = R_V^* - R_S^*$ achieved by the receiver under voluntary disclosure.

\begin{figure}[htbp]
    \centering
    \includegraphics[width=0.6\textwidth]{regression_decision_sweep.png}
    \caption{Regression Strategic Advantage $\Delta^*(E)$ as a function of feature-label correlation $\rho_{XY}$ and feature-preference correlation $\rho_{XU}$ (under channel noise $q = 0.3$, employing a coolwarm colormap). The dashed black contour delineates the boundary of the sufficient condition (Theorem 3.4).}
    \label{fig:regression_decision_sweep}
\end{figure}
\end{sourceblock}

\begin{faqblock}
\faqitem{Process}{What parameter sweep is executed to analyze strategic advantage?}
\faqanswers
\faqans{Process}{Sweeps continuous parameters ($\rho_{XY}$, $\rho_{XU}$) to map regions of positive strategic advantage $\Delta^*$.}
\end{faqblock}

\subsubsection{Regression Decision Sweep}

\begin{sourceblock}{Paragraph 6 - Regression Sweep Results}
Figure~\ref{fig:regression_decision_sweep} displays the strategic advantage heatmap for the regression game under a fixed channel noise $q = 0.3$. We sweep $\rho_{XY}$ and $\rho_{XU}$ on a $25 \times 25$ grid. 
\begin{itemize}
    \item \textbf{Observations:} A positive strategic advantage (indicated by red regions) is attained when the feature-label correlation $\rho_{XY}$ and the feature-preference correlation $\rho_{XU}$ have the same sign. In this regime, the sender's disclosure incentive aligns with the receiver's predictive benefit: the sender voluntarily reveals features when they are large in the direction of the label, which coincides with the instances where features are highly predictive. Conversely, when the correlations have opposite signs, voluntary disclosure degrades accuracy compared to OLS Omit baselines (blue regions).
    \item \textbf{Sufficient Condition Validation:} The dashed black contour shows the boundary where the sufficient condition (Theorem~\ref{thm:sufficient_condition}) holds. As predicted, this region is a subset of the true positive strategic advantage region, confirming that the condition is sufficient but conservative due to the Cauchy-Schwarz approximation of the preference bias.
\end{itemize}
\end{sourceblock}

\begin{faqblock}
\faqitem{Consequence}{When does voluntary disclosure outperform mandatory OLS baseline?}
\faqitem{Conclusion}{How does the sufficient condition boundary compare to the empirical results?}
\faqanswers
\faqans{Consequence}{When $\rho_{XY}$ and $\rho_{XU}$ have matching signs, aligning the sender's disclosure motive with receiver accuracy.}
\faqans{Conclusion}{Sufficient condition contour (Theorem 3.5) lies strictly inside the positive advantage region, validating its sufficiency and conservativeness.}
\end{faqblock}

\begin{sourceblock}{Figure 2 - Classification Results}
\begin{figure}[htbp]
    \centering
    \includegraphics[width=0.48\textwidth]{classification_piecewise_linear.png}
    \hfill
    \includegraphics[width=0.48\textwidth]{classification_decision_sweep.png}
    \caption{Classification decision results. Left: Piecewise-linear envelope of strategic risks for different defaults $i$, validating Theorem~\ref{thm:piecewise_linear_noise}; Right: Heatmap of strategic advantage $\Delta^*(q)$ over noise $q$ and alignment noise $\alpha$ (consistent coolwarm colormap).}
    \label{fig:classification_decision_results}
\end{figure}
\end{sourceblock}

\begin{faqblock}
\faqitem{Process}{What classification features are illustrated in Figure 2?}
\faqanswers
\faqans{Process}{Left: Piecewise-linear risk envelope over noise $q$; Right: Strategic advantage heatmap over noise $q$ and alignment noise $\alpha$.}
\end{faqblock}

\subsubsection{Classification Decision Sweep}

\begin{sourceblock}{Paragraph 7 - Classification Sweep Results}
Figure~\ref{fig:classification_decision_results} (left) plots the strategic risks $R^i_S(q)$ for each of the $K=3$ default classes alongside the optimal strategic risk envelope $R^*_S(q) = \min_i R^i_S(q)$, for a fixed alignment noise $\alpha = 0.15$.
\begin{itemize}
    \item \textbf{Piecewise-Linear Structure:} In accordance with Lemma~\ref{lem:class_strategic_risk}, each individual risk line is affine in $q$. The lower envelope $R^*_S(q)$ is concave and piecewise-linear, partitioned into three linear segments separated by transition points near $q \approx 0.35$ and $q \approx 0.70$. This confirms Theorem~\ref{thm:piecewise_linear_noise}: the optimal default class changes from class 2 to class 1, and finally to class 3 as the channel noise increases, partitioning the noise interval.
    \item \textbf{Advantage Heatmap:} Figure~\ref{fig:classification_decision_results} (right) sweeps the channel noise $q$ and alignment noise $\alpha$. The strategic advantage is largest when the alignment noise $\alpha$ is low (indicating high alignment between preference $U$ and label $Y$) and the channel noise $q$ is small (ensuring features are successfully transmitted). As the alignment noise $\alpha$ increases beyond $0.5$, the preference becomes decoupled from the target class, and the strategic advantage vanishes.
\end{itemize}
\end{sourceblock}

\begin{faqblock}
\faqitem{Consequence}{How is Theorem 4.5 validated by the piecewise-linear structure?}
\faqitem{Consequence}{What environmental factors maximize strategic advantage in classification?}
\faqanswers
\faqans{Consequence}{Lower envelope of risk lines forms a concave piecewise-linear curve with distinct segments separating optimal defaults.}
\faqans{Consequence}{Low alignment noise $\alpha$ (high alignment) and low channel noise $q$ (successful transmission).}
\end{faqblock}

\subsection{Offline Generalization Gaps}
\label{sec:exp_offline}

\begin{sourceblock}{Paragraph 8 - Offline Learning Gaps Overview and Fig 3}
Next, we evaluate the offline learning curves by performing multiple experiments with different training sample sizes $n \in \{50, 100, 200, 500, 1000, 2000, 4000\}$. We explicitly track the excess strategic risk $R_S(\hat{b}_n) - R_S(b^*)$ (for regression) and $R_S(\hat{\Theta}_n) - R_S(\Theta^*)$ (for multiclass classification) for each distinct training size $n$, represented as discrete points in our scatter plots, and averaged over $15$ independent runs.

\begin{figure}[htbp]
    \centering
    \includegraphics[width=0.48\textwidth]{regression_offline_learning.png}
    \hfill
    \includegraphics[width=0.48\textwidth]{classification_offline_learning.png}
    \caption{Offline learning scatter plots showing the decay of the excess strategic risk on a log-log scale for varying training sizes $n$. Left: Regression excess risk converging at a rate matching the $\mathcal{O}(d^3/\sqrt{n})$ bound; Right: Multiclass classification excess risk converging at a rate matching the $\mathcal{O}(\sqrt{Kd/n})$ Natarajan dimension bound.}
    \label{fig:offline_learning_curves}
\end{figure}
\end{sourceblock}

\begin{faqblock}
\faqitem{Process}{What offline training sizes $n$ and independent runs are evaluated?}
\faqanswers
\faqans{Process}{$n \in \{50, 100, 200, 500, 1000, 2000, 4000\}$ over 15 independent runs, plotting average excess risk.}
\end{faqblock}

\begin{sourceblock}{Paragraph 9 - Offline Generalization Gaps Results}
\begin{itemize}
    \item \textbf{Regression Convergence:} Figure~\ref{fig:offline_learning_curves} (left) shows the convergence rate of the strategic empirical risk minimizer (ERM) on a log-log scale. The excess risk decays rapidly, demonstrating that a sample size of $n \ge 1000$ is sufficient to achieve an excess strategic prediction risk of less than $0.05$. The decay matches the theoretical $\mathcal{O}(n^{-0.5})$ generalization rate proved in Theorem~\ref{thm:excess_risk_guarantee}.
    \item \textbf{Classification Convergence:} Figure~\ref{fig:offline_learning_curves} (right) validates the sample complexity of the joint selection algorithm in classification. The excess risk converges to zero as $n$ approaches $4000$, validating Lemma~\ref{lem:simultaneous_convergence} and Theorem~\ref{thm:classification_excess_risk}. The convergence rate matches the sublinear rate $\mathcal{O}(n^{-0.5})$ governed by the Natarajan dimension of multiclass linear boundaries.
\end{itemize}
\end{sourceblock}

\begin{faqblock}
\faqitem{Consequence}{What empirical convergence rates are observed for regression and classification?}
\faqanswers
\faqans{Consequence}{Both show sublinear $\mathcal{O}(n^{-0.5})$ decay of excess risk, validating the theoretical sample complexities.}
\end{faqblock}

\subsection{Online Regret Convergence}
\label{sec:exp_online}

\begin{sourceblock}{Paragraph 10 - Online Regret Overview and Fig 4}
Finally, we analyze the online sequential learning setting under bandit feedback, running the algorithms over $T = 3000$ steps.

\begin{figure}[htbp]
    \centering
    \includegraphics[width=0.48\textwidth]{regression_online_regret.png}
    \hfill
    \includegraphics[width=0.48\textwidth]{classification_online_regret.png}
    \caption{Online expected cumulative regret scatter plots under bandit feedback. Left: Regression regret under bandit OGD using one-point gradient estimates for $d=1$ and $d=3$, showing $\mathcal{O}(d\sqrt{t})$ sublinear growth with dimension-dependent reference benchmarks; Right: Classification regret under EXP3+OGD for $K=2$ and $K=3$, confirming the sublinear $\mathcal{O}((K \ln K)^{1/3} t^{2/3})$ regret bound with class-dependent benchmarks.}
    \label{fig:online_regret_curves}
\end{figure}
\end{sourceblock}

\begin{faqblock}
\faqitem{Process}{What step horizon $T$ is used in the online sequential simulations?}
\faqanswers
\faqans{Process}{$T=3000$ steps, evaluated under bandit feedback where unrevealed information is hidden.}
\end{faqblock}

\begin{sourceblock}{Paragraph 11 - Online Regret Results}
\begin{itemize}
    \item \textbf{Regression Bandit OGD:} Figure~\ref{fig:online_regret_curves} (left) displays a scatter plot of the cumulative regret of the one-point bandit gradient estimator. For both dimensions $d=1$ and $d=3$, the regret grows sublinearly, tightly bounded by their respective dimension-dependent $\mathcal{O}(d\sqrt{t})$ theoretical benchmarks. Comparing the two, we clearly observe that the higher dimension $d=3$ incurs proportionally larger cumulative regret. This is due to the $d/\delta_t$ scaling in the one-point spherical gradient estimate, matching Theorem~\ref{thm:bandit_feedback_rate}.
    \item \textbf{Classification EXP3 + OGD:} Figure~\ref{fig:online_regret_curves} (right) presents a scatter plot of the expected cumulative regret of our decoupled EXP3+OGD algorithm (Algorithm 1) for $K=2$ and $K=3$ classes. The regret grows sublinearly, conforming to the distinct $K$-dependent $\mathcal{O}((K \ln K)^{1/3} t^{2/3})$ reference curves proved in Theorem~\ref{thm:classification_regret_bound}. The curve for $K=3$ classes exhibits a slower convergence rate and higher regret than $K=2$ classes, as the meta-algorithm must explore a larger discrete choice space (EXP3 arms) to identify the optimal default class.
\end{itemize}
\end{sourceblock}

\begin{faqblock}
\faqitem{Consequence}{How does regression regret scale with dimensionality?}
\faqitem{Consequence}{How does classification regret scale with class size $K$?}
\faqanswers
\faqans{Consequence}{Regret scales with dimension $d$, matching the $\mathcal{O}(d\sqrt{t})$ bound from bandit gradient estimation.}
\faqans{Consequence}{Incurs higher regret for larger $K$, matching the discrete arm exploration bounds of $\mathcal{O}((K \ln K)^{1/3} t^{2/3})$.}
\end{faqblock}
